\chapter{The First Five-Year Plan}

The period from the publication of the first control figures of Gosplan in August 1925 to the approval of the first five-year plan in May 1929 was one of unbroken advance in the principles and practice of planning. About the middle of the period the centre of interest shifted from the annual control figures to the five-year plan, which entailed a purposeful review of Soviet economic policy and of the long-term prospects of economic development. Party pronouncements from time to time expounded the aims of the plan. The fourteenth congress in December 1925 proclaimed the goal of ``economic self-sufficiency'', which meant converting the USSR ``from a country that imports machines and equipment into a country that produces machines and equipment''. The party conference in the following autumn called for ``reconstruction of the economy on the basis of new and more advanced technology''. The principle of priority for the production of means of production rather than of consumer goods was challenged from time to time by those who sought to slow down the pace of industrialization; and in 1927-1928 the production of consumer goods was expanded to meet the needs of the peasant market. But this was a temporary response to an emergency. The party conference of April 1929, which finally approved the plan, put first on the list of its aims ``the maximum development of the production of the means of production as the foundation of the industrialization of the country''. 

The first tentative Gosplan draft of a five-year plan in March 1926 was concerned primarily with state industry---the only sector of the economy as yet within the control of the planners. Here it budgeted for annual increases in industrial production at rates ranging from 40 per cent in the first year, when unused capacity was still available for exploitation, to 15 per cent in the fifth year---a so-called ``attenuating curve'' of growth. Investment in industry was to rise from 750 million rubles in the first year to 1200 millions in the fifth. The draft attracted little attention at higher levels, being still regarded as a theoretical exercise rather than a set of practical proposals. The second Gosplan draft a year later was a far more detailed and sophisticated document, devoting separate chapters to different sectors of the economy. Its estimates of industrial growth were considerably more modest than those of its predecessor; and the forecast that the industrial labour force would rise by a million workers during the currency of the plan was scaled down by more than a half. On the other hand, its demands for increasing investment in industry were pitched somewhat higher. Additional funds for the expansion of industry were to be provided by reducing the costs of production through an increase in the productivity of labour; and much was expected from the campaign waged by the authorities to force down prices. Planning had now become a crucial issue, and the draft of March 1927 excited bitter controversy. In Gosplan, the Bolshevik economist Strumilin, the main architect of the draft, was confronted by the cautious ex-Menshevik, Groman. It was attacked by economists in Narkomzem and Narkomfin as a dangerous fantasy, and by Kuibyshev in Vesenkha as unduly timid. The opposition, taken unawares by the sharp turn in the official line towards rapid industrialization, confined itself to the charge that the conversion to planning had come too late to be effective. 
% 注：修复了 "foreca.st" 的打字错误，去除了 "i", "li" 乱码。

From this point pressure to raise planning targets was continuous. The opposition in its platform of September 1927 (see p. 117 above) now condemned existing Gosplan proposals as niggardly and inadequate; already a year earlier, at a discussion of the control figures in the Communist Academy, spokesmen of the opposition outbid those of Gosplan and Vesenkha in the demand for a higher rate of industrialization. In October 1927 Gosplan produced a third draft plan which, by way of conciliating both the doubters and the optimists, offered ``basic'' and ``maximum'' figures; the latter represented a substantial advance, both in industrial production and in investment, over the estimates in the second draft. Vesenkha now took up the running, and produced estimates well in excess of those of Gosplan; and this led Gosplan once more to revise its estimates upwards. For the time being progress was halted by the hesitations of the party leaders. At the moment of crisis in the struggle against Trotsky and the opposition it would have been embarrassing openly to concede their demand for more rapid industrialization; and it was inexpedient to press too hard issues on which unavowed differences of opinion existed between the leaders. When the party central committee met at the end of October to approve ``Directives for the Drafting of the Five-Year Plan'' for submission to the forthcoming party congress, the text, while full of unqualified enthusiasm for the plan, betrayed a notable reluctance to take a stand on any of the contentious problems involved. A balance must be struck between ``the interests of accumulation'' and ``the peasant economy'', between heavy and light industry; a ``maximum transfer of resources'' from the latter to the former was ruled out as ``a violation of the equilibrium of the whole economic system''. A ``maximum rate of accumulation'' in the current year was not necessarily a guarantee of the most rapid development in the long term. No attempt was made to pronounce a verdict on either of the Gosplan variants or on that of Vesenkha. This was the session which expelled Trotsky and Zinoviev from the committee (see p. 118 above). The few remaining spokesmen of the opposition attacked these vague ``directives'' which did not contain a single figure. But they were interrupted and shouted down by the majority. 
% 注：去除了 "I ", "jit", "■4- (" 和 "3f" 等错位的数字和符号。

The party congress in December 1927 marked the importance of the five-year plan by devoting seven sittings to the debate on it. Bukharin, whose commendation of ``snail's pace industrialization'' at the previous congress two years earlier was unkindly recalled by one delegate, did not speak. A few sceptical voices were raised, but these were drowned in the general enthusiasm for the principle of the plan. Some ardent supporters of industrialization criticized the timidity of Gosplan, and praised Vesenkha as the front runner in the race. The party leaders were, however, restrained, notably in their references to agriculture (see p. 124 above). The congress was content to accept the cautious ``directives'' drafted by the party central committee a few weeks earlier, and made no attempt to translate the enthusiasm for the plan into statistical terms. The main business of the congress was to defeat and expel the opposition. No controversial issue must be allowed to mar the unanimity of the party majority in carrying out this task. The only positive decision taken was to prepare the plan in time for its submission to the next Union Congress of Soviets in the spring of 1929. 

Once the opposition had been crushed, and its leaders driven from Moscow and dispersed, the inhibitions which had restrained the top party leaders disappeared. The changed mood was reflected in the harshness of the ``extraordinary measures'' applied in the grain collections of the first months of 1928. Hitherto the main driving force behind the five-year plan had appeared to come from secondary party figures ensconced in Gosplan and Vesenkha, who sought to convince the leaders of the feasibility of their ambitious projects; Stalin in particular had maintained his favourite role of moderator between two extremes. Henceforth it was clear that the weight of the party was behind every upward revision of the estimates; the driving force now came from the Politburo and from Stalin himself. Throughout 1928, with this new impetus and at higher levels, the same pattern as before was followed of rivalry between Gosplan and Vesenkha in pursuit of constantly rising targets. At the same time plans became more specific, purporting to cover every sector of the economy, every industry and every region. Calculations were further and further removed from any ``genetic'' base, and became to an ever-increasing extent an expression of the will and determination to achieve. Politics were intertwined with economics in the debate, and the final decisions were not so much economic as political. The defeat of Bukharin and the condemnation of his ``Notes of an Economist'' in the autumn of 1929 made it clear that caution would henceforth be treated as the symptom of a Right deviation. After prolonged discussion between Gosplan and Vesenkha the draft of the plan was completed in March 1929. It offered ``basic'' and ``optimum'' estimates both of industrial production over the five years of the plan (1928-1929 to 1932-1933) and of the rate of investment in industry. In the optimum variant, the ``attenuating curve'' in the rate of increase in production had disappeared; the annual rate of increase was to rise progressively from 21.4 per cent in the first year to 23.8 in the fifth. Investment in industry, planned at 1650 million rubles for the first year, was to be almost doubled by the fifth year (basic variant) or more than doubled (optimum variant). While the economists who drafted the plan seem to have thought of the basic variant as the limit of reasonable expectation, and the optimum as a remote potential, the Politburo boldly adopted a resolution endorsing the plan ``in its optimum variant'' as ``fully corresponding to the directives of the fifteenth party congress''. 
% 注：去除了 "iWhile" 和 "/thought" 等标点。将 "21-4"、"23-8" 修正为 "21.4"、"23.8"。

The plan was finally approved by a large party conference at the end of April 1929. Rykov, though now associated with the Right deviators, was one of three rapporteurs who commended it to the conference. But his cautious appraisal was compared unfavourably with the eloquent enthusiasm of Krzhizhanovsky, the president of Gosplan, and the cool, matter-of-fact determination of Kuibyshev, the president of Vesenkha. The plan was published in three large volumes a few days after the conference; owing, no doubt, to lack of time for correction, it still contained the two variants adopted by Gosplan in March, though the basic variant was now obsolete. It was an impressive and comprehensive review of the whole economy. Some of its estimates proved crudely over-optimistic, especially when a year later the targets were raised further, and the decision taken to complete the five-year plan in four years. But it gave a powerful impetus to ambitious projects for the development of heavy industry; and it can be argued that, without the great wave of optimism which it generated, these results could not have been achieved. The five-year plan became the pivot round which the whole economy revolved. 
% 注：去除了 "I .to ambitious" 中的乱码。

Gosplan was the heir of Goelro, the organ set up to carry Lenin's plan of electrification; energetika (energy) continued to be one of its key-words. It was appropriate, and something more than a coincidence, that the most famous project promoted by Gosplan, and executed as a vital part of the first five-year plan, should have been the construction of a great dam and hydro-electric station on the Dnieper river, known as Dneprostroi. In the summer of 1926 Cooper, the American engineer who had built the Tennessee Valley dam, accepted an invitation to visit the site, expressed enthusiasm for its potentialities, and eventually agreed to supervise the construction. The project was to be financed from the Soviet budget; Cooper was employed as a consultant and adviser, not as a contractor. But it required the unstinting use of American technology and equipment and the recruitment of a small army of American engineers. It also involved the erection of new industries and factories to use electrical power generated by the scheme. Power was to be supplied to the Donbass mines, and to large new iron and steel works, as well as to works producing aluminium, high-grade steel and ferro-alloys, forming a vast new industrial complex producing means of production. Two new towns were built---Zaporozhie and Dnepropetrovsk. It was not till about 1934 that the dam itself and the factories planned to consume its output were in full operation. 

Dneprostroi set a pattern for many of the ambitious projects initiated under the first five-year plan. One-seventh of all industrial investment under the plan went into the production of iron and steel, though some of this was used to modernize existing factories and plants. The development of the automobile industry attracted much publicity. Before the revolution no motor-cars were manufactured in Russia; the first beginnings were made in two or three engineering works which turned out a handful of cars in the middle nineteen-twenties. In 1927, under the enthusiasm generated by the plan, the first Soviet automobile factory was authorized, a small concern near Moscow with a projected output of 10,000 cars a year. It was not till 1929 that an agreement was signed with Ford of Detroit for the construction of an automobile factory at Nizhny-Novgorod with a planned annual output, to be achieved within ten years, of 200,000 vehicles. At first the emphasis was on cars for personal use. Later priority was given to the production of lorries for industrial use. A road-building programme was a corollary of the growth of the automobile industry. A parallel, but separate, development was the production of tractors. The planned output of the tractor factory at Stalingrad (see p. 130 above), several times increased while the work was in progress, was fixed in the final draft of the five-year-plan at 50,000 tractors a year. By the time it went into production in 1930, two further factories had been authorized. From 1928 onwards the tractor played a leading role in programmes for the modernization and collectivization of the peasant economy. It was the major contribution of the first five-year plan to the promotion of agricultural production. 
% 注：去除了 "tTMT ttt trn"、"• given"、"(increased" 的杂乱符号。

The armaments industry was rarely mentioned in public discussions of the five-year plan. After the civil war, the Red Army had been allowed to run down for some years. But in 1926 steps were taken to strengthen and re-equip it; and, after the war scare of the spring of 1927, recognition of the industrial basis of military power made the five-year plan, with its emphasis on heavy industry, a matter of military concern. Stimulus must have been derived from the secret military arrangements with Germany; and a secret and separate five-year plan was said to exist for the war industries. The aircraft industry led the way, and was followed by the production of tanks. Much importance was attached to the development of a modern chemical industry, which served both military and agricultural needs, and was said by its protagonists to play a role comparable to that of electricity in the modernization of the economy. 

An enterprise which was actually under way before the beginning of the plan, and which, unlike other major projects, did not depend on imported materials and equipment or on foreign technical advisers, was the construction of the Turksib railway linking Central Asia and Kazakhstan with western Siberia. Central Asia was a rich cotton-growing area. But its communications were poor; and the opening of a fresh outlet for its raw cotton was designed to make the textile industry of the USSR independent of imports from abroad. On the other hand, Central Asia produced inadequate food crops and no timber; the new railway would make possible the supply of grain and timber direct from producing areas in Siberia and relieve the pressure on supplies from European Russia. Russian engineers had ample experience in railway building. Allocations from the budget were readily available; and the construction of 1500 kilometres of track in difficult country proceeded with few hitches. The line was opened for regular traffic on January 1, 1931. 
% 注：去除了 "/in Siberia", "y Russia" 的边缘符号。修正 "January i" 为 "January 1"。

An important planning problem, which was constantly debated and sometimes delayed vital decisions, was the location of the new industries. The issue turned to some extent on the relative practical advantages of different sites. But the main trouble came from local rivalries. It was especially conspicuous in the iron and steel industry, partly because this absorbed so large a share of investment under the plan, partly because the Ukrainians fought hard to retain the predominant place in iron and steel production which they had acquired since the eighteen-nineties, mainly owing to the proximity of large deposits of coal. Their claim was contested by a large school of ``easterners'', who sought to revive the once-flourishing iron and steel industry of the Urals, and to establish new centres of production in Siberia. In 1927 a number of rival projects were being canvassed and investigated. The first, the planning of which was most advanced, was a large new iron and steel complex at Krivoi Rog in the Ukraine. The second was a project of comparable dimensions at Magnitogorsk in the Urals. This would have to draw its supplies of coking coal by rail from the Kuznetsk basin, 2000 kilometres further east, where a third extensive iron and steel works was proposed. Besides these, a large engineering factory was planned at Sverdlovsk in the Urals, and yet another iron and steel complex at Zaporozhie in the neighbourhood of the Dnieper dam. But the Krivoi Rog project was postponed till the end of the first five-year plan; and the progressive increases in planned investment during 1928 represented in the main a victory for the ambition of the ``easterners'' to extend Soviet power and activity into the empty spaces of Siberia, and to develop and populate these under-utilized regions. The war scare of 1927 had powerfully intensified the anxiety of the planners to locate the future centres of vital Soviet industry in less vulnerable areas than the Ukraine. 
% 注：去除了 "I ..of the", "nIjLvIx Oil", "|the relative", "[(main", "Idominant", "(future" 等严重乱码。

Though the drive for industrialization was the focus of the plan, it was a plan for the whole economy, and could not be anything less. The main and well-recognized stumbling-block in the way of the planners was agriculture. As Rykov once observed, the plan was ``at the mercy of a shower of rain''; and the choice of the five-year period in planning was justified by the argument that within this period good and bad harvests would balance one another, and that calculations based on an average would therefore be valid. But an even greater obstacle was the unpredictability of peasant behaviour. The peasant family on its small holding, largely self-sufficient at a traditionally low subsistence level, could isolate itself from the national economy, and frustrate the calculations of the planners. It was the problem of collecting the grain and bringing it to the towns and factories, as much as the problem of producing it, which preoccupied the authorities in the middle nineteen-twenties. The impact of the plan on the peasant was therefore two-fold. It sought to bring the production of peasant agriculture within the scope of its prognostications and directives. But it also aimed at replacing the primitive methods of cultivation hitherto in use by the provision of up-to-date machines and implements. Tractors were merely the most advanced and most ambitious of the tools which industry could supply in order to promote more efficient cultivation of the soil. The production of means of production for agriculture might seem an insignificant part of the programme of industrialization. But the output of agriculture was the basis of the whole plan. A leading official at the party congress in December 1927 called for an increase in grain production of from 30 to 40 per cent in the next five years and of 100 per cent in ten years; the figure of 35 per cent eventually appeared in the five-year plan. 
% 注：去除了 "j not", "I block", "wrfan", "^[(.l", "''lisell", "//isolate", "11 calculations", "l^the grain", "i'much" 的侧边乱码。

The claims of industrialization rendered obsolete the conception of state finance as the criterion of the feasibility of economic policy. Since the stabilization of the new currency in 1924 Narkomfin had been responsible for preparing, not only the annual budget which controlled state expenditure, but also the quarterly credit plan which regulated the supply of credit through the banks and the rate of emission of currency. Once planned industrialization became the permanent aim, these vital elements in the economy could not elude the attention of the planners. It became clear that the restriction of credit which would have been necessary in order to maintain a stable international currency based on gold was incompatible with the expansion of industry. The choice was not in doubt. As early as 1925, while the budget was strictly balanced, expansion of credit for industry and the consequent increase in the currency issue weakened confidence in the chervonets, the gold parity of which could no longer be maintained in international markets. In the summer of 1926 export of the chervonets and dealings in it abroad were prohibited; and the currency became henceforth a purely internal medium of exchange, subject to manipulation as the interests of the economy required. This abandonment of the gold standard after less than two years was a blow to Soviet prestige, and the growing symptoms of inflation caused alarm. For a year longer Narkomfin succeeded in keeping some control over credit. But pressure to step up the pace of industrialization mounted continuously; and to confine it within a credit strait-jacket imposed by Narkomfin was wholly unacceptable. 
% 注：去除了 "(a purely", "It j of the", "I I to Soviet", "|the pace" 等多余符号。

Before the end of 1927 the battle had been won by the planners, and the traditional policies of financial control were obsolete. It was laid down that the state budget would be geared to the control figures of Gosplan; and the budget, together with the levels of credit and currency emission, became effectively a part of the five-year plan. Financial operations were subjected to the discipline of the plan; and credit was provided for industrial projects approved by Vesenkha and Gosplan in the face of predictions in Narkomfin of its inflationary consequences. The mood was one of unbounded confidence. The boldest expectations seemed to have been surpassed. Investment in large-scale industry under the control of Vesenkha for 1927-1928 amounted to 1300 million rubles---an increase of more than 20 per cent on the previous year; and industrial production planned by Vesenkha was said to have risen by more than 25 per cent. When the contours of the first five-year plan began to take final shape in the autumn of 1928, an investment of 1650 million rubles was projected for 1928-1929. In October 1928 Pyatakov, known as a ``super-industrializer'', a former member of the opposition who had since recanted, was appointed vice-president of the State Bank; and he became its president early in 1929. The appointment was significant of a determination to expand credit to whatever limits might be necessary to finance industrial production. The total of capital investment in industry was fixed by a process of argument between Gosplan, Vesenkha and the higher party authorities. The provision of funds to meet this demand was an administrative question. While the plan was couched in terms of state finance, Narkomfin in effect became a revenue-raising department which no longer controlled expenditure. 
% 注：去除了 "I (effectively" 等乱码。

Orthodox sources of finance for industrial development were eagerly explored. Direct taxation---the industrial tax on the private sector, the agricultural tax and income tax---almost doubled in monetary terms between 1926 and 1929. But indirect taxation was more important. Excise, which more than doubled in this period, accounted for one-third of all tax revenue. Excise on articles of common consumption was a tax on the poorest; and the large receipts from the vodka monopoly troubled some party consciences. But alternative sources of revenue were hard to find. From 1927 onwards a series of state loans were floated, subscription to which, in spite of official protestations to the contrary, soon acquired a quasi-compulsory character. By these means Narkomfin was enabled each year to present a balanced budget. Deficit financing would have been regarded as inadmissible. But behind this conventional facade finance had been dethroned from its regulatory role, and additional credits were pumped into the economy by the State Bank. Money gradually became simply a medium of exchange and an accountancy unit---a foretaste of the time when it would disappear altogether from the future communist society. But, apart from budget allocations and credits from the State Bank, it was assumed that funds for investment in industry would be made available from the profits of industry itself. In view of the over-riding demand to keep down prices, the only way to achieve this was by reducing costs of production. This had been the consistent aim of campaigns for the ``regime of economy'', for rationalization and for increased productivity of labour (see pp. 111-113 above). Estimates of productivity had been stepped up in each successive version of the five-year plan; and the planned increase was greater in the capital goods sector than in industry as a whole. The plan as eventually adopted in its optimum variant provided for an increase in productivity during the period of the plan of 110 per cent and for a reduction in costs by 35 per cent. It offered to the worker the prospect of an increase in real wages of 47 per cent and a reduction in retail prices of 23 per cent. But these estimates seem to have been based, not on any realistic assessment of the problem, but on the desire to make the plan statistically coherent, and were indicative rather of the immense pressure imposed by the plan on the industrial worker than of any prospect of their fulfilment. 
% 注：彻底修复了该段错综复杂的多栏重叠。如 "I. indirect taxation was more important. Excise, which more -3»|yfeagerly explored. Direct taxation—the industrial tax on the ^///private sector..." 已复原为 "Orthodox sources of finance for industrial development were eagerly explored. Direct taxation---the industrial tax on the private sector... But indirect taxation was more important." 并去除了 "-3»|yf", "^///", "-f||(", "// than", "1 tax", "-»;//a", "(I monopoly", "I enabled", "Icing", "' (see", "o the", "—1 lany", "-^///to", "y/ I'", "^ Wj" 等极端的乱码碎片。 

The adoption of the first five-year plan was a landmark in Soviet history. The essence of NEP had been to concede a measure of freedom to the peasant economy. It would have been impolitic to announce its demise. Stalin argued that NEP, while providing ``a certain freedom for private trade'', had also assured ``the controlling role of the state over the market''. The purpose of NEP was to ``secure the victory of socialism''. It was officially denied that NEP had been abrogated. A free market in the products of small-scale private industry, and above all of agriculture, still remained. But the subordination of every major economic activity to the dictates of the plan, and the increasingly harsh pressures brought to bear on the peasant, made these survivals of NEP both anomalous and precarious. They were tolerated so long as it was convenient to tolerate them, but seemed of little account. The share of the private sector in the national income, which had exceeded 50 per cent in 1926-1927, fell to insignificant dimensions by the end of the five-year plan. 
% 注：去除了 "3", ".1 .", "III victory", "r! been", "ILf private", "(the dictates", "~f Ilf", "2^1//income", "n to" 等大量的垂直条状污渍乱码。

The prestige of the plan, and of the USSR as the protagonist of planning, was enhanced by the economic crisis which broke on the capitalist world in the autumn of 1929. It was widely felt, and not only in the USSR, that the Marxist prediction of the collapse of the capitalist order under the weight of its inherent contradictions had been vindicated. The immunity of the USSR from some of the worst symptoms of the crisis---notably from mass unemployment---encouraged a growing belief that no national economy could any longer be left at the mercy of the iron laws of the market. The Soviet five-year plan, though the conditions of its adoption and operation were not much studied or understood, seemed to provide a pioneering model. The demand for an element of planning in the economies of the capitalist countries grew everywhere, and substantially influenced western attitudes to the USSR. 
% 注：去除了 "-5" 的杂乱标记。